% 本文件由 LaTeX 排版 Agent 根据有效 translation.json 自动生成。
% author=John Chrysostom; work_id=2305; title=论得撒洛尼后书讲道集

\chapter{\WorkTitle{论得撒洛尼后书讲道集}}

% structure: p0001 source=h1 target=omitted reason=page-title-already-rendered-by-parent

讲道一\newline{}讲道二\newline{}讲道三\newline{}讲道四\newline{}讲道五\par

\SourceRecord{Translated by John A. Broadus.；Nicene and Post-Nicene Fathers, First Series；Vol. 13.；Edited by Philip Schaff.；Buffalo, NY: Christian Literature Publishing Co.,；1889.；<http://www.newadvent.org/fathers/2305.htm>.}

% structure: p0001 source=h1 target=section reason=page-title

\section{得撒洛尼后书讲道 第一篇}

\Emph{引言}\par

他在前一封书信中说：\Quote{我们日夜祈祷，渴望见到你们}，并且\BibleQuote{我派了弟茂德前去}\BibleRef{得前 3:1}，通过这些话语，他显明了想要来到他们中间的渴望。因此，当他或许没有时间前往，并圆满他们信德中所欠缺的部分时，出于这个原因，他又加上了第二封书信，以他的写作来弥补他不在场的缺失。因为他没有离开，我们可以由此推测：因为他在此书信中说：\BibleQuote{我们因主耶稣基督的来临恳求你们}\BibleRef{得后 2:1}。因为在他第一封书信中，他说：\BibleQuote{关于时期和季节，你们不需要有人写给你们}\BibleRef{得前 5:1}。所以，如果他去了，就没有必要写作了。但由于这问题被推迟了，故而他添加了这封书信，正如他在致弟茂德的书信中说：\BibleQuote{他们败坏了一些人的信德，说复活已经过去}\BibleRef{弟后 2:18}；为使信友们从此对任何伟大或辉煌之事不再抱希望，而可能在他们所受的苦难中软弱下来。\par

因为那希望支持了他们，不让他们屈服于当前的邪恶，魔鬼切望剪除这希望，视之为一种锚，当他不能说服他们未来的事是虚假时，便另寻途径，收买了某些有害的人，试图欺骗那些相信的人，使他们认为那些伟大而辉煌的事已经应验。因此，这些人当时说复活已经过去。但现在他们说审判和基督的来临近了，为的是甚至使基督也陷入谎言中，并且指出此后没有报应、没有审判台、也没有惩罚和报复给那些曾对他们行恶的人，他们既可以令这些行恶者更加胆大，又可以使那些受迫害者更加沮丧。而且，比这一切更糟的是，有些人仅仅试图传扬仿佛由保禄所说的话语，而另一些人甚至伪造由他所写的书信。为此，他堵住了他们的一切进路，说：\BibleQuote{你们不要迅速动摇思想，也不要惊慌，无论因灵、因言语或因如同由我们所发出的书信}\BibleRef{得后 2:2}。\Quote{也不要因灵}，他指着假先知说。那么，我们要如何分辨他们呢？他说。正因这个缘故，他加上：\BibleQuote{我保禄亲笔问安，这是每封书信的记号：我这样写。愿我们主耶稣基督的恩宠与你们众人同在}\BibleRef{得后 3:17}。他这里的意思并不是说这就是记号——因为很可能别人也模仿了这个——而是说我亲笔写这问安，如在我们中间现在的习惯一样。因为通过署名，发信人的作品便被知晓。但他安慰他们，因为他们被烦恼过分折磨；既从他们当前的状况称赞他们，又从未来的远景、惩罚以及为他们预备的善报鼓励他们；并且更清楚地详述了这主题，不是揭示时间本身，而是指明时间的记号，即假基督。因为软弱的灵魂最得确证，不是仅仅听，而是得知更具体的事。\par

基督也极其关注这一点，坐在山上时，祂极其详细地向门徒谈论了终局。为什么呢？为的是不让那些引进假基督和伪基督的人有机可乘。祂自己也给出了许多记号，其中一个，也是最重要的，说：\BibleQuote{福音必须传遍万民}\BibleRef{玛 24:14}，另一个是，他们不应被祂的来临所欺骗。\BibleQuote{如同闪电}\BibleRef{玛 24:27}，祂说，祂将要来临；不隐藏在任何角落，而是照耀各处。它不需要任何人来指明，它将如此辉煌，甚至闪电也不需要人来指明。祂也在某处论到假基督，当祂说：\BibleQuote{我因我父的名而来，你们却不接待我；若另有人因自己的名而来，你们倒要接待他}\BibleRef{若 5:43}。祂还说，那些接连不断的不可言说的灾难是它的记号，而且厄里亚必须来。\par

得撒洛尼人当时确实困惑，但他们的困惑对我们却有益处。因为不仅为他们，而且为我们，这些事也是有益的，好使我们摆脱幼稚的寓言和老妇人的愚妄。你们不是常常在孩童时听到人们大谈特谈假基督的名字和他屈膝的事吗？因为魔鬼将这些事撒在我们还年幼的心思中，好让这教训随着我们成长，使他能欺骗我们。因此，保禄在谈论假基督时，如果这些事是有益的，他本不会忽略它们。所以我们不要探究这些事。因为他不会这样屈膝而来，而是\BibleQuote{相反一切称为神或受人敬拜者，以致他坐在天主的殿中，自充为天主}\BibleRef{得后 2:4}。因为正如魔鬼因骄傲而堕落，同样，那被他所感化的人也被膏以骄傲。\par

因此，我恳求你们，让我们都热切地远离这种情感，免得我们陷入他的定罪，免得我们承受同样的惩罚，免得我们分受那报复。他说：\Quote{不可选新奉教的人，免得他自高自大，落在魔鬼的定罪里。}\BibleRef{弟前 3:6}所以自高自大的人，与魔鬼同受惩罚。因为\Quote{骄傲的开端，就是不认识上主。}\BibleRef{德 10:12}骄傲是罪的开端，是趋向恶的第一冲动与动向。也许它确实是根源与基础。因为\Quote{开端}的意思，要么是趋向恶的第一冲动，要么是基础。如同有人说，贞洁的开端是避见不当之物，那是第一冲动；但如果我们说，贞洁的开端是禁食，那是基础与确立。同样，骄傲是罪的开端。因为每一样罪都始于它，并由它维持。因为无论我们做什么善事，这一恶习都不容它们存留、不容它们不坠落，反而像某种根，不让它们安稳不动，这从以下事例明显可见：看法利塞人做了什么，但对他毫无益处。因为他没有拔除那根，那根反而败坏了他所有的行为，因为根留下来了。从骄傲生出藐视穷人、贪求财富、喜爱权势、渴望大荣耀。这样的人容易报复侮辱。因为骄傲的人，即使受上级侮辱也不能忍受，何况下级？但不能忍受侮辱的人，也不能忍受任何损害。看，骄傲是如何成为罪的开端。\par

但不认识主如何是骄傲的开端呢？理当如此。因为凡按应当的方式认识天主的人，凡认识天主子如此谦卑自抑的人，就不会自高自大。但不认识这些事的人，就会自高自大。因为骄傲膏抹他，使他傲慢。请告诉我，那些向教会发动战争的人，他们说认识天主，这是从何而来？岂不是来自傲慢？看，不认识主把他们投入何等悬崖！因为天主喜爱痛悔的心灵\BibleRef{咏 50:17}，另一方面他却\Quote{拒绝骄傲的人，赐恩宠给谦逊的人。}\BibleRef{伯前 5:5}因此，没有比骄傲更大的邪恶了。它使人变成魔鬼，傲慢、亵渎、背信，并使人渴望死亡与凶杀。骄傲的人总是生活在苦恼中，总是生气，总是不快乐。没有什么能使他饕足。如果他看见君王向他俯身下拜，他并不满足，反而更加激怒。因为如同爱钱的人，收取得越多，就越想要；同样，骄傲的人，享受的荣耀越多，就越渴望。因为他们的情欲增长了；因为这是一种情欲，情欲不知限度，只在它杀死拥有者时才停止。你岂不见醉汉总是口渴吗？因为这是一种情欲，不是本性的渴望，而是某种扭曲的疾病。你岂不见那些患\OriginalTerm{病态饥饿}{bulimy}的人总是饿吗？因为这是一种情欲，正如医家子弟所说，已经超越了本性的界限。那些好管闲事、过分好奇的人，无论学到什么，都不停止。因为这是一种情欲，没有限度。\BibleRef{德 23:17}再者，那些喜爱淫乱的人，也不能停止。经上说：\Quote{对于行淫者，所有的饼都是甜的。}他不停止，直到被吞噬。因为这是一种情欲。\par

但它们确实是情欲，却不是不可医治的，而是可医治的，甚至比身体的病痛更容易医治。因为我们若愿意，就能熄灭它们。那么人如何能熄灭骄傲呢？藉着认识天主。因为骄傲源于不认识天主，如果我们认识祂，一切骄傲就都被驱逐了。想想地狱。想想那些比你优秀得多的人。想想你的罪。想想你该受天主惩罚的事有多少。你若想到这些，就能很快降伏你骄傲的心，很快使它谦卑。但你不能做这些事吗？你太软弱了吗？思虑当下之事，思虑人性本身，思虑人的虚无！当你看见一具尸体被抬过市集，孤儿们跟在后面，寡妇捶胸痛哭，仆人们哀悼，朋友们神色沮丧，请反思当下之事的虚无，以及它们与影子或梦境并无不同。\par

这对你不适用吗？想想那些非常富有的人，他们无论如何死于战争；环顾那些属于显贵名流的房屋，如今已夷为平地。思考他们曾如何强大，而今连他们的纪念也不复存在。因为，你愿意的话，每天都能发现这类例子——统治者的更替——富人财产的充公。因为\Quote{许多暴君曾坐在地上，而那从未被想到的人却戴上了王冠。}\BibleRef{德 11:15}这些事不是每天发生吗？我们的事不是像旋转的轮子吗？你若愿意，既读我们自己的书，也读外教人的书，因为其中也充满这样的例子。你若因骄傲而轻视我们自己的书，你若钦佩哲学家的著作，就去他们那里吧。他们会教导你，讲述古代的灾难，诗人、演说家、诡辩家以及所有历史学家也是如此。从各方面，你若愿意，都能找到例子。\par

若你们不愿接受这些事，就反思我们自己的本性：它由什么构成，又归于何处。试想，当你睡眠时，你又有何价值？连一个小野兽岂不能毁灭你吗？因为常常有一只小动物从屋顶掉下，使许多人失明，或造成其他危险。但什么？你不比所有野兽更软弱吗？但你说什么？说你以理性见长？看，你并没有理性：因为骄傲是缺乏理性的标志。告诉我，你究竟为何高傲？是因你身体良好的构造吗？但此胜利的奖赏属于无理性的受造物；强盗、凶手和盗墓者也拥有它。但你以你的智慧为傲吗？骄傲并不是智慧的证明。因此，你首先使自己失去了成为明达之人的机会。让我们降低高傲的思想。让我们节制、谦卑和温良。因为就连基督也如此称他们为有福的，超越一切，说：\BibleQuote{神贫的人是有福的。}\BibleRef{玛 5:3} 同样，祂高声说：\BibleQuote{你们跟我学吧，因为我是良善心谦的。}\BibleRef{玛 11:29} 为此，祂洗了门徒的脚，给我们留下了谦卑的榜样。让我们从这一切中获得益处，使我们能够获得为爱祂的人所预许的福分，藉着祂的恩宠和慈爱，等等。\par

\SourceRecord{Translated by John A. Broadus.；Nicene and Post-Nicene Fathers, First Series；Vol. 13.；Edited by Philip Schaff.；Buffalo, NY: Christian Literature Publishing Co.,；1889.；<http://www.newadvent.org/fathers/23051.htm>.}

% structure: p0001 source=h1 target=section reason=page-title

\section{《得撒洛尼后书》讲道词 第二篇}

\BibleRef{得后 1:1-2}\par

\BibleQuote{\Person{保禄}、\Person{息耳瓦诺}与\Person{弟茂德}，致书给在天主我们的父及主耶稣基督内的得撒洛尼人的教会。愿恩宠与平安，由天主父及主耶稣基督赐与你们。}\par

世人多半以一切言行讨好像和高于自己的人的欢心；他们视此为大事，认为自己幸福，若能达成此目的。但是，如果得人的欢心是如此巨大的好处，那么得天主的欢心是何等重大呢？因此，他总这样开始他的书信，并祈求这恩宠临于他们，因为他知道，若这一恩宠得着，以后便无任何可悲之事，无论何种困难，都将消解。为使你明白这点，请想\Person{若瑟}：他年轻时是奴隶，没有经验，尚未成长，突然一家之管治交在他手中，他须向一位埃及主人交代。你们知道那民族多么易怒、不宽恕，况且加上权柄和能力，他们的怒气更大，被权力所燃。这一点也从后来发生的事明显可见。当主母控告时，他忍耐了。然而，那不是持有衣服之人的事，而是被剥去衣服之人遭受暴行。因为他本应说：若他听见你如你所说提高了声音，他必已逃走；若他有罪，他必不等主人到来。但他丝毫不考虑这类事，却不合理地完全被怒气所制，将他投入监牢。他是如此轻率之人。而且，即便从其他事上，他也可推断此人的良好品格和聪慧。然而，因他极不合理，他从未考虑任何这类事。因此，那位与如此严厉的主人相处、并被托付管理全家的异乡人、孤单无经验之人，当天主倾注丰盛恩宠于他时，他度过了一切，仿佛他的试探从未存在过：主母的诬告、死亡的危险、监狱，最终到了王座。\par

所以这位有福之人明见天主的恩宠何其大，因此他将此恩宠祈求赐予他们。此外，他也成就另一件事：希望使他们对书信剩余部分怀有好感，即使他责备和规劝他们，他们也不至于脱离他。为此，他在一切之先提醒他们天主的恩宠，软化他们的心，这样，即使有苦难，但想到他们藉此恩宠得救脱离更大的恶，他们便不至于因较小的磨难而绝望，反而能从中获得安慰。正如他在另一封信中曾说过：\Quote{若我们还在为仇敌的时候，曾藉着祂圣子的死与天主和好了，那么，在和好之后，更必因祂的生命而得救了。} \BibleRef{罗 5:10}\par

\BibleQuote{愿恩宠与平安，}他说，\BibleQuote{由天主父及主耶稣基督赐与你们。}\par

第三节 \BibleQuote{我们常该为你们感谢天主，弟兄们，这真是相称的。} \BibleRef{得后 1:3}\par

这又是极大谦逊的标记。因为他引导他们反省并思考：若我们的善行不是先由人称赞，而是由天主称赞，那么我们更当如此。在其他方面，他也鼓舞他们的精神，因为他们所受的苦并不值得流泪哀叹，而是值得向天主感恩。但是，如果\Person{保禄}为别人的善行感恩，那么那些不但不感恩，反而因此嫉妒的人，将受何等的惩罚呢？\par

\BibleQuote{因为你们的信德格外增长，你们众人彼此相爱的心也格外增多。}\par

\BibleQuote{因为你们的信德格外增长，你们众人彼此相爱的心也格外增多。}\par

因为特别爱某一个人有什么益处呢？那是人的爱。但如果它不是人的爱，而是为天主的缘故去爱，那么就要爱所有人。因为天主甚至命令我们要爱仇人。如果祂命令我们爱仇人，何况那些从未得罪过我们的人呢？但是，你说：\Quote{但是，我并没有做这些事。}但当一个人被诽谤，而你不堵住说话者的口，不相信他的话，不阻止他时，这能证明什么样的爱呢？他说：\Quote{而且你们众人彼此相爱的心也越发增长。}\par

第4节。\Quote{甚至我们在天主的各教会里为你们夸口。}\newline{}\BibleRef{得后 1:4}\par

的确，在前一封书信中，他说马其顿和阿哈雅的各教会都传扬开了，听说了他们的信德。他说：\Quote{所以我们不需要说什么，因为他们自己传报我们怎样进到你们那里。}\newline{}\BibleRef{得前 1:8} 但在这里，他说：\Quote{所以我们夸口。}那么，这话是什么意思呢？在那里，他说他们不需要他的教训，但在这里，他没有说我们教训他们，而是说\Quote{我们夸口}，以你们为荣。因此，如果我们既为你们感谢天主，又向人夸口，你们更应当为自己的善行如此行。因为如果你们的善行值得别人夸耀，它们怎么会值得你们哀叹呢？这说不通。他说：\Quote{甚至我们在天主的各教会里为你们的信德和忍耐夸口。}\par

在这里，他表明已经过了许多时间。因为忍耐是通过长时间显明出来的，不是两三天。而且他不仅仅说忍耐。忍耐严格来说是指在尚未享受所应许的福分时的表现。但这里他指的是更大的忍耐。那是什么样的忍耐呢？就是在迫害中所显明的忍耐。他说：\Quote{因你们的信德和忍耐，在一切迫害和你们所忍受的患难中。}因为他们与敌人生活在一起，这些敌人不断从各方面试图伤害他们，而他们表现出坚定不可动摇的忍耐。让所有为了人的庇护而转投其他教义的人脸红吧。因为当传道刚刚开始，靠每日工资生活的穷人承受来自统治者和国家首脑的敌意——当时没有一位国王或总督是信徒——并且进行着不可调和的战争，即便如此，他们仍不动摇。\par

第5节。\Quote{这正是天主公义审判的明证。}\newline{}\BibleRef{得后 1:5}\par

看他如何为他们寻求安慰。他说，我们感谢天主，我们在人前夸口：这些确实尊荣。但一个受苦之人最渴望的是脱离苦难，并对恶待他们的人施以报复。因为灵魂软弱时，最追求这些；有哲学精神的灵魂甚至不追求这些。那么他为什么说\Quote{公义审判的明证}呢？这里他提到双方——作恶的和受苦的——的报应，仿佛在说：天主的公义得以彰显，祂既给你加冕，也惩罚他们。同时，他也安慰他们，表明他们因自己的劳苦和辛劳得冠冕，按公义的比例。但他先提到他们的部分。因为即使一个人强烈渴望报复，他首先渴望的是奖赏。因此他说：\par

\Quote{使你们配得天主的国，你们正是为这国受苦。}\par

这种事之所以发生，不是因为伤害他们的人比他们更有能力，而是因为必须这样他们才能进入天主的国。他说：\Quote{我们必须经历许多艰难，才能进入天主的国。}\newline{}\BibleRef{宗 14:22}\par

第6、7节。\Quote{因为天主既是公义的，就必将患难报应那加患难给你们的人，并给你们这受患难的人，与我们同得平安，在主耶稣从天上同祂大能的天使显现的时候。}\par

这里的\Quote{既是}意思是\Quote{因为}，我们也在谈论非常明确、不可否认的事情时使用它，代替说：\Quote{因为是极其公义的。}他说：\Quote{天主既是公义的}，就必惩罚这些人。就好像他说：\Quote{如果天主关心人间事务，} \Quote{如果天主眷顾。}他不凭自己的意见说话，而是把它当作公认的事实；就像有人说：\Quote{如果天主恨恶恶人，}以此迫使他们承认天主确实恨恶他们。因为这类说法是最无可争议的，既然他们自己也认识到那是公义的。因为如果在人看来这是公义的，那么在天主看来更加如此。\par

他说：\Quote{将患难报应那加患难给你们的人，并给你们……平安。}那么，报应是相等的吗？绝非如此。请看他如何通过下文表明报应更严厉，而\Quote{平安}也大得多。这也是另一个安慰：他们在患难中有同伴，也在报应中有同伴。他使他们与那些做了无限更多、更大工作的人同得冠冕。然后他加上时期，并通过描述将他们的思想提升向上，仿佛已经用言语打开天堂，将其呈现在他们眼前；他将天使大军围绕祂，从地方和侍从两方面扩大画面，使他们得到些许安慰。他说：\Quote{给你们这受患难的人，与我们同得平安，在主耶稣从天上同祂大能的天使显现的时候。}\par

第8节。\Quote{在火焰中，报应那不认识天主和那不听从我们主耶稣福音的人。}\newline{}\BibleRef{得后 1:8}\par

如果那些没有服从福音的人遭受惩罚，那么那些不仅不服从而且还折磨你们的人，岂不更要受罚吗？请看他的智慧：他此处没有说\Quote{那些折磨你们的人}，而是说\Quote{那些不服从的人}。这样，即使不是为你们的缘故，但为祂自己的缘故，也必须惩罚他们。这话是为了确证，他们必须受罚是全然必要的；之前的话则是为了使他们也受到尊荣，因为他们为你们的缘故而受苦。前者使他们相信惩罚，后者使他们欢喜，因为他们为了自己所受的而受苦。\par

这些话是对他们说的，但也适用于我们。因此，当我们受苦时，让我们思想这些事。让我们不要因他人受罚而欢喜，仿佛报了仇，而要看作我们自己逃脱了这样的惩罚和报应。他人受罚对我们有何益处呢？我恳求你们，不要有这样的心。让我们因天国的盼望而被引向德性。因为那极其有德的人，既不受恐惧驱使，也不被天国的盼望所动，唯独为基督，如同\Person{保禄}那样。然而，让我们即使如此思想天国的福乐、地狱的痛苦，从而规范并教训自己；让我们以此方式引领自己走向当实践的事。当你看到今生任何美好而伟大的事物时，思想天国，你便会视之为无物。当你看到任何可怕的事物时，思想地狱，你便会嘲笑它。当你被肉欲占据时，思想那火，也思想罪本身的快乐，它毫无价值，甚至没有快乐。如果此地所立的法律的恐惧有如此大的力量，能使我们远离恶行，那么对将来之事的记忆，那不朽的报复，那永久的惩罚，其力量岂不更大？如果对地上君王的恐惧能使我们远离如此多的恶事，那么对永恒之王的恐惧岂不更大？\par

那么，我们如何能常常持有这种恐惧呢？如果我们不断聆听圣经。因为只要看见一具死尸就足以使心神消沉，更何况地狱和不灭的火，更何况那不死之虫。如果我们常常思想地狱，我们就不会轻易堕入地狱。为此，天主警告了惩罚；如果思想它没有大益处，天主就不会警告它。但因为记忆它能产生大善，故此祂将它的恐怖放入我们灵魂，如同有益的良药。因此，我们不可忽视由此而来的大益处，而应常常提及它，在午餐时，在晚餐时。因为谈论愉快之事对灵魂毫无益处，反而使它更松懈；而谈论痛苦悲伤之事，则能斩断灵魂中一切松弛放荡，转化它，使它在软弱时刚强起来。那谈论剧场和演员的人，不仅无益于灵魂，反而更点燃它，使它更粗心。那关心并忙碌于他人事务的人，常常因这好奇心使灵魂陷入危险。但那谈论地狱的人，不冒任何危险，并使灵魂更加清醒。\par

但你害怕这些话的冒犯性吗？那么，你若缄默，就能熄灭地狱吗？你若谈论，就能点燃它吗？无论你谈论与否，火焰都在喷涌。让它被常常提起，使你不要堕入其中。灵魂若为地狱忧虑，就不可能轻易犯罪。请听这至善的劝诫：经上说：\BibleQuote{你要记得你的末路，你的结局}\BibleRef{德 28:6}，你就永远不会犯罪。惧怕交账的灵魂，不能不迟于过犯。因为恐惧在灵魂中强劲有力，不容世上任何事物存留其中。如果关于地狱的谈论如此使灵魂谦卑降伏，那么这思念常驻灵魂，岂不比任何火更能净化它吗？\par

让我们不要那么记念天国，而要记念地狱。因为恐惧比应许更有力。我知道许多人会轻看万种福乐，只要能免除惩罚；因为如今对我来说，逃脱报应、不受罚已经足够了。没有一个将地狱放在眼前的人会堕入地狱。没有一个轻视地狱的人能逃脱地狱。因为正如我们中间那些惧怕审判座的人不会落入他们手中，而那些轻视他们的人正是那些落在他们之下的；这里也是如此。如果\Place{尼尼微人}没有惧怕毁灭，他们就会被倾覆；但因他们惧怕，就没有被倾覆。如果在\Person{诺厄}时代他们惧怕洪水，他们就不会被淹没。如果\Place{索多玛人}惧怕，他们就不会被焚烧。轻视警告是极大的恶事。轻视警告的人，很快就会在警告的执行中体验到它的真实。没有什么比谈论地狱更有益了。它使我们的灵魂比任何银子更纯净。请听先知说：\BibleQuote{你的判断常在我面前}\BibleRef{咏 16:22}。因为虽然它使听者痛苦，但大大有益于他。\par

因为确实，凡是有益的事都是如此。就连药物和食物，起初也会令病人不适，随后才带来益处。如果我们不能忍受严厉的言辞，显然我们也无法承受实际的苦难。若无人能忍受关于地狱的谈论，那么显然，若迫害来临，没有人能在烈火和刀剑面前站稳。让我们操练我们的耳朵，不要过于娇嫩脆弱；因为这样，我们才能逐渐承受那些实际的事物。如果我们习惯于听闻可怕的事，我们也会习惯于承受可怕的事。但如果我们如此软弱，连言语都无法忍受，那么当面对实际时，我们怎能站立得住？你看见有福的保禄如何鄙视世上的一切，连一个接一个的危险都不当作考验吗？为什么呢？因为他已习惯于为了天主的旨意而藐视地狱。为了基督的爱，他甚至将经历地狱视为无物；而我们却连为了自己的益处也不愿忍受关于地狱的谈论。现在，你们稍微听了一下，就请回去吧；但我恳求你们，如果你们有丝毫的爱德，要经常回到关于这些事的谈论中来。它们不会伤害你们，即使不能使你们受益，也必定会对你们有益。因为我们的灵魂是根据我们的言论而被塑造的。经上说：\Quote{邪恶的交谈败坏善良的风俗。}因此，美好的交谈也能改善灵魂；同样，令人畏惧的谈论能使灵魂清醒。因为灵魂就如同一块蜡。如果你用冷淡的言论去涂抹，就会使它变硬、变得麻木；但如果你用火热的言论，就会使它融化；融化之后，你就能按自己的意愿塑造它，并在其上刻印君王的肖像。因此，让我们堵住耳朵，不听那些虚妄的言论。这并非小恶，因为一切邪恶都由此而生。\par

如果我们的心思习惯于专注于神圣的谈论，它就不会转向其他谈论；而不转向其他谈论，也就不会陷入恶行。因为言语是行为的途径。首先我们思想，然后我们说话，然后我们行动。许多人，原本清醒，却往往因可耻的言语进而行可耻的事。因为我们的灵魂并非生来为善或为恶，而是因选择而成为二者之一。因此，如同船帆将船带往风所吹的方向，更确切地说，如同舵在顺风时引导船只，同样，如果善言从有利的方向吹来，思想便会在无危险中航行。但若风向相反，它们甚至会淹没理性。因为话语之于灵魂，就如同风之于船只。你愿意，就能移动并转动它。为此，一位劝诫者说：\Quote{愿你的全部谈论都在至高者的法律中。}（\BibleRef{德 20:20}）因此，我劝你们，当我们从乳母手中接过孩子时，不要让他们习惯于老妇的传说，而要让他们从幼年起就学习审判和惩罚的存在；要把这恐惧铭刻在他们心中。这恐惧一旦扎根，就会产生极大的益处。因为一个从幼年起就学会被这种期待所抑制的灵魂，不会轻易摆脱这种恐惧。它如同一匹驯服的马，服从缰绳，地狱的思想如同骑手坐在其上，使之步履有序，这样它就会说出有益的言语；无论是青春、财富、孤儿身份，还是其他任何事物，都不能伤害它，因为它的理性如此坚定，能够抵御一切。\par

让我们用这些谈论来规范我们自己，也规范我们的妻子、仆人、儿女、朋友，如果可能，甚至我们的仇敌。因为藉着这些谈论，我们能够斩断大部分的罪恶，并且宁可谈论令人忧愁的事，也不谈论令人愉快的事，这一点从以下事例显而易见。请告诉我：如果你进入一个举行婚宴的家庭，一时你会为景象所喜悦，但随后离开时，你会因自己未能拥有同样多的财富而痛苦。但如果你进入哀悼者的家，即使他们非常富有，当你离开时，你会感到更加轻松。因为在那里，你没有产生嫉妒，反而因你的贫穷得到了安慰和慰藉。你从事实中看到：财富并非美善，贫穷并非邪恶，它们都是无关紧要的事。同样，现在，如果你谈论奢侈，你会更加折磨你的灵魂，因为它也许无法奢侈。但如果你指责奢侈，并引入关于地狱的谈论，那件事会使你振奋，并产生许多喜乐。因为当你想到奢侈丝毫不能保护我们免于那烈火时，你就不会追求它；但如果你反思奢侈甚至更容易点燃那烈火，你不仅不会追求，反而会避开并拒绝它。\par

我们不要回避谈论地狱，好使我们能避免地狱。不要驱除对惩罚的记忆，好使我们能逃脱惩罚。如果那个富翁曾默想那火，他就不至于犯罪；但因他从未记念它，所以他跌入其中。告诉我，人啊，你即将站在基督的审判台前，你却谈论一切事，而不谈论那事？当你在一位置法官前有一案件，常常只关乎言语，你昼夜不停，无时无刻不谈别的事，总是谈论那案件；但当你即将为你的全部生命交账，并接受审判时，你甚至不能容忍别人提醒你那审判吗？因此，一切都败坏和荒废，因为当我们即将为今生的事站在人间的法庭前时，我们调动一切，恳求众人，常为此焦虑，做一切事为它；但当我们不久后要来到基督的审判台前时，我们却什么也不做，既不借自己，也不借他人；我们不恳求审判主。然而祂赐给我们长久的容忍时期，不将我们正中罪恶时夺去，却允许我们脱去它们，那良善与慈爱并未留下任何属于祂自己的事不做。但一切都无济于事；因此惩罚将更重。但愿天主不许如此！因此，我恳求你们，让我们即使现在也要警惕。让我们把地狱放在眼前。让我们思考那不容情的账目，这样，思考这些事，我们既能避免罪恶，也能选择美德，并能获得那应许给爱祂之人的福分，借着恩宠和慈爱，等等。\par

\SourceRecord{Translated by John A. Broadus.；Nicene and Post-Nicene Fathers, First Series；Vol. 13.；Edited by Philip Schaff.；Buffalo, NY: Christian Literature Publishing Co.,；1889.；<http://www.newadvent.org/fathers/23052.htm>.}

% structure: p0001 source=h1 target=section reason=page-title

\section{《得撒洛尼人后书讲道辞第三篇》}

\Emph{\BibleRef{得后 1:9-10}}\par

\Emph{\Quote{他们要受惩罚，就是永远丧亡，与主的面，和祂威能的光荣隔绝，当祂来临时，要在祂的圣人身上受光荣，在一切信者身上显奇妙。}}\par

有许多人不是靠远离罪恶而形成良好的希望，而是认为地狱并不像所说的那样可怕，比所威胁的更温和，是暂时的，而非永远的；他们对此大加思辨。但我能从许多理由，并从关于地狱的那些说法本身证明，它不仅是更不温和，而且比所威胁的更可怕得多。但我现在不打算讨论这些事。因为即使只是言语的恐惧就足够了，尽管我们不完全展开其含义。但要证明它不是暂时的，请听保禄现在论及那些不认识天主和不信从福音的人时说：\Quote{他们要受惩罚，就是永远丧亡。}那么，暂时的东西如何能是永久的呢？他说：\Quote{与主的面。}这是什么意思？他在这里想说这是多么容易的事。因为那时他们极其傲慢，他说，不需要多费周折；只要天主来临并被看见，所有人都陷入惩罚和报应。祂的来临对某些人确实是光，但对另一些人却是报应。\par

\Quote{和祂威能的光荣隔绝，}他说，\Quote{当祂来临时，要在祂的圣人身上受光荣，在一切信者身上显奇妙。}\par

天主受光荣吗？是的，他说，在众圣人身上。如何？因为当那些极其傲慢的人看见那些曾受他们鞭打、轻蔑、嘲笑的人，如今竟亲近祂时，这就是祂的光荣，或者更确切地说，是他们和祂的光荣；祂的光荣在于祂没有舍弃他们；他们的光荣在于他们配得了如此大的尊荣。正如拥有忠信的人是祂的财富，同样，拥有那些将享受祂恩惠的人也是祂的光荣。良善者的光荣在于拥有那些可以施与他恩惠的人。\Quote{并显奇妙，}他说，\Quote{在一切信者身上，}也就是说，\Quote{通过信者。}请再看，这里的\Quote{在}被用于\Quote{通过。}因为通过他们，祂被显示为可钦佩的，当祂使那些可怜、不幸、遭受无数苦难并相信了的人达到如此辉煌时，祂的大能便显现出来；因为尽管他们似乎在世上被抛弃，然而在那里他们仍享受大光荣；那时特别显示出天主所有的光荣和大能。如何？\par

\Quote{因为我们在那日子对你们的作证被信从了。}\par

第11节。\Quote{因此，我们也常为你们祈祷。} \BibleRef{得后 1:11}\par

这就是说，当那些曾遭受无数苦难、被诱使背弃信仰而并未屈服、反而相信了的人被公开展示时，天主受光荣。那时也显示出这些人的光荣。\Quote{未死以前，不要称任何人是有福的，}经上说。\BibleRef{德 11:28}为此他说：在那一日将显明那些信了的人。\Quote{因此，我们也常为你们祈祷，}他说，\Quote{愿我们的天主使你们堪当蒙召，并以大能成全你们各种良善的意愿和信德的工作。}\par

\Quote{愿祂使你们堪当，}他说，\Quote{蒙召；}因为他们并未蒙召。因此他又加上：\Quote{并成全各种良善的意愿。}因为那个身穿污秽衣服的人也蒙了召，但他没有坚守他的召叫，反而因此更被抛弃。\Quote{蒙召，}即被召到婚宴。因为五个童女也蒙了召。\Quote{起来，}经上说，\Quote{新郎来了。} \BibleRef{玛 25:6}她们也准备了，但未能进去。但祂说的是另一种召叫。因此，为表明祂所说的是哪种召叫，祂又加上：\Quote{并用大能成全各种良善的意愿和信德的工作。}这就是那召叫，祂说，我们所寻求的。请看祂多么温和地贬抑他们。为免他们因过分的称赞而变得虚荣，好像已经做了大事，并且变得懈怠，祂指出他们在今生仍有不足之处。这也就是祂在致希伯来人书中所说的。\Quote{你们与罪恶争斗，还没有抵抗到流血的地步。} \BibleRef{希 12:4} \Quote{直到完全中悦，}他说，即祂的喜悦、说服、完全保证。这就是说，愿天主的说服得以实现，你们毫无欠缺，成为祂所愿意的那样。\Quote{并用大能成就各种信德的工作，}他说。这是什么？就是忍受迫害的坚韧，他说，使我们不致丧气。\par

第12节。\Quote{好使我们的主耶稣基督的名，在你们身上受光荣，你们也在祂身上受光荣，都是按照我们的天主和主耶稣基督的恩宠。} \BibleRef{得后 1:12}\par

祂在那里论及荣耀，在这里也论及它。祂说，他们被荣耀，以至他们甚至可以夸耀。祂说，更为甚者，他们同时也荣耀天主。祂说，他们将接受那荣耀。但在这里，祂也意指：由于主被荣耀，仆人们也被荣耀。因为那些荣耀他们主人的人，他们自己更被荣耀，既因那件事本身，又因它之外的事。为基督的缘故而受困苦便是荣耀，祂到处称此事为荣耀。我们越是忍受任何可耻的事，我们就越显赫。然后又表明这件事本身也出于天主，祂说：\Quote{按照我们的天主和主耶稣基督的恩宠}；也就是说，祂已亲自将这恩宠赐给我们，为叫祂在我们内被荣耀，也叫我们在祂内被荣耀。祂如何在我们内被荣耀？因为我们不把任何事放在祂之前。我们如何在祂内被荣耀？因为我们从祂领受了能力，使我们完全不屈服于加在我们身上的邪恶。因为当试炼发生时，天主同时被荣耀，我们也如此。因为他们荣耀祂，因祂如此坚固了我们；他们钦佩我们，因我们使自己成为配得的。而这一切都是藉着天主的恩宠成就的。\par

第二章 第一至二节。\BibleQuote{现在我们请求你们，弟兄们，关于我们的主耶稣基督的来临，和我们的聚集到祂跟前；为使你们不要很快地被摇动你们的心。}\par

复活何时发生，他没有说，但他说了不会在现在。\Quote{和我们的聚集到祂跟前。} 这也不是小事。请看，劝勉又如何再次伴随着称赞和鼓励，因为祂和所有圣人必将与我们一同显现。他在这里论述复活和我们的聚集。因为这些事将同时发生。他振奋他们的心神。\Quote{为使你们不要很快地被摇动，} 他说，\Quote{也不要被困扰，无论是因神、因言语，或因一封似乎发自我们的书信，以为主的日子现已临到。}\par

他在这里似乎是暗示，有些人到处传扬一封伪造的书信，仿佛出自\Person{保禄}，并展示此信，说主的日子近了，从而引诱许多人陷入错误。因此，为免他们受骗，\Person{保禄}藉他所写的内容提供保障，并说：\Quote{不要被困扰，无论是因神或因言语}。他所说的意思是：即使有人拥有预言之神而这样说，也不要相信。因为当我和你们在一起时，我已告诉你们这些事，所以你们不应改变从所受教导而来的确信。或者如此理解：\Quote{因神}，他这样称呼假先知，他们藉不洁之神说他们所说的话。因为这些人更愿被人相信，不仅试图用说服人的言语来欺骗（因为这是他指出的，说\Quote{或因言语}），他们还展示一封假冒发自\Person{保禄}的信，宣告同样的事。因此，为也指出这一点，他补充说：\Quote{或一封看似出自我们的书信。} 在各方面稳妥了他们之后，他这样阐述自己的教导，并说：\par

第三节至第四节。\BibleQuote{不要让任何人以任何方式欺骗你们：因为那日子不会来到，除非那背教先发生，那罪恶之人，丧亡之子，被揭露出来，他敌对并高举自己超越一切称为神或受人敬拜的；以致他坐在天主的殿中，自充为神。}\par

他在这里论述\OriginalTerm{假基督}{Antichrist}，并揭示伟大的奥秘。什么是\Quote{背教}？他称之为\OriginalTerm{背教}{the falling away}，因他将要毁灭许多人，使他们背教。所以祂说：\Quote{若是可能，连被选者也要被迷惑。}\BibleRef{玛 24:24} 他又称他为\OriginalTerm{罪恶之人}{the man of sin}。因为他要做无数恶事，并促使别人去做。但他称他为\OriginalTerm{丧亡之子}{the son of perdition}，因为他也要被毁灭。但他是谁？是撒殚吗？绝不是，而是一个人，允许撒殚在他内完全运行。因为他是一个人。\Quote{并高举自己超越一切称为神或受人敬拜的。} 因为他不会引入偶像崇拜，而是一种反对天主的人；他将废除所有神明，并命令人敬拜他自己以代替天主，他将坐在天主的殿中，不单是在耶路撒冷的那殿，也在每一个教会中。\Quote{自充为天主，} 他说；他没有说他宣告，而是他努力显示。因为他将行大奇事，并显奇妙征兆。\par

第五节。\BibleQuote{你们不记得吗？当我还在你们那里时，我已把这些事告诉你们了？}\BibleRef{得后 2:5}\par

你看到有必要持续说同样的事，并用同样的话详述它们吗？因为看哪，他们亲耳听到他说这些事，却又需要被提醒。因为正如他们听到关于苦难的事，他说：\Quote{因为确实，当我们在你们那里时，我们预先告诉你们，我们将要受苦难，}\BibleRef{得前 3:4} 但他们仍然忘记了，他又用书信坚固他们；同样，听了关于基督的来临，他们又需要书信来使他们安顿。因此他提醒他们，表明他并没有说什么新奇的事，而是他常说的。\par

因为如同农夫的情形，种子确实一次被撒在地里，却不会恒久存留，还需要许多准备工作；如果他们不松土、不覆盖所撒的种子，他们就是为聚集的雀鸟播种。同样，我们若不以不断的记忆覆盖所撒下的，我们就只是把它撒在空中了。因为魔鬼将它夺去，我们的懈怠毁掉它，太阳晒干它，雨冲走它，荆棘窒息它。所以，一次撒种后离开是不够的，还需要许多恒心，驱赶雀鸟，拔除荆棘，用大量的土填满石地，抑制、围护并除去一切有害的东西。但在土地的情形中，一切都靠农夫，因为它是一个无生命的主题，只准备被动。而在精神的土壤中则完全不同。一切不都是教师的部分，而是至少一半，甚至更多，属于门徒。我们的部分确实是撒种，但你们要做那些为提醒你们而说的话，通过你们的行为显明果实，连根拔起荆棘。\par

因为财富实在是荆棘，不结果实，既不美观，又不悦用，伤害那些触摸它的人，不仅自己不结果实，甚至还阻碍那正要生长的。财富就是如此。它不仅不结永恒的果实，反而阻碍那些想要获得它的人。荆棘是无理性的骆驼的食物；它们被吞吃，被火焚烧，毫无用处。财富也是如此，毫无用处，只用来点燃火炉，照亮那如烤炉燃烧的日子，喂养无理性的情欲、报复和愤怒。因为那吃荆棘的骆驼也是如此。那些熟悉此事的人说，没有一种动物像骆驼那样难以和解、那样闷闷不乐和报复。财富就是如此。它滋养灵魂中无理性的情欲，却刺伤理性的部分，如同荆棘的情形。这种植物坚硬粗糙，而且自己生长。\par

让我们看它是如何生长的，好将它连根拔除。它生长在陡峭、多石、干燥、没有水分的地方。因此，当一个人粗糙而陡峭，即不慈悲时，荆棘就在他里面生长。但当农夫的儿子们想要拔除它们时，他们不用铁器。那么怎样呢？他们放火烧掉，用这种方式除去土地的一切坏品质。因为仅仅砍掉上部，而根留在下面不够；甚至根除根也不够（因为它因坏品质留在地里，如同某种瘟疫侵袭身体后，仍留下残余），从上而来的火，吸走荆棘的所有水分，如同毒液，通过热量从地心抽出。如同放在患处的拔火罐将所有紊乱吸向自己，火也抽出荆棘中一切恶劣品质，使土地洁净。\par

那么我为什么说这些呢？因为你们应当清除对财富的一切贪恋。我们也有一种火能抽出灵魂的这种恶劣品质；我指的是圣神之火。如果我们让它作用于它们，我们不仅能干枯荆棘，也能干枯它们的水分，因为如果它们扎得很深，一切就都白费了。请看，有一个富人进来了，或者一个富女人？她不关心如何听天主的圣言，而是如何炫耀，如何以排场就座，如何以许多荣耀，如何以衣着的昂贵超越所有别的女人，并通过服饰、表情和步态使自己更显尊贵。她所有关心和顾虑是：那个女人看见我了吗？她羡慕我吗？我的美丽是否恰当地展现？这样她的衣服不会腐烂，也不会撕裂；所有心思都在于此。同样，富人进来，意图在穷人面前展示自己，以他周身的衣服和众多奴隶震慑他们，这些奴隶也站在周围，驱赶人群。但他因极大的骄傲甚至不屑于这样做，反而认为这是一件如此不配绅士的举动，尽管极度自负，却不能忍受去做，而是把它交给奴隶。因为做这事确实需要真正奴仆和厚颜的举止。然后当他坐下时，家中的烦扰立刻侵入，四面八方地分散他的注意力。占据他灵魂的骄傲满溢出来。他认为他是在向我们也向民众，甚至向天主施恩，因为他进入了天主的殿。但如此被炙烤的人，怎能被治愈呢？\par

请告诉我，如果有人去医生的诊所，不向医生求恩惠，反而认为是在给医生施恩，并且拒绝为他的伤口请求药物，反而关心他的衣服；他会得到任何益处而离开吗？我想不会。但是，请允许我告诉你们这一切的原因。他们认为当他们进入这里时，他们是进入我们面前；他们认为他们听到的是从我们听到的。他们不放在心上，不考虑他们正在进入天主面前，是祂在向他们说话。因为当读经员站起来说：\Quote{主这样说，}执事站着命令所有人肃静时，他这样说并不是为了尊敬读经员，而是为了尊敬那通过他向所有人说话的那一位。如果他们知道是天主通过祂的先知说这些话，他们就会抛弃所有骄傲。因为当统治者对他们说话时，他们尚且不让心思游荡，何况当天主说话时呢？我们是仆人，亲爱的。我们不讲自己的事，而是讲天主的事；从天上来的信件每天被宣读。\par

求你们告诉我，我恳求你们，如果现在，当我们全体在场时，有一个人走进来，腰束金带，挺直身子，带着一副了不起的神情说，他是地上君王差来的，并且带来了致全城关于重要事务的书信；难道你们不都转向他吗？难道你们不无需执事下令就保持深沉的静默吗？我确实这样想。因为我常在这里听到君王们的书信被宣读。那么，如果有人从君王那里来，你们全都留意；而先知从天主而来，从天上说话，却没有人留意吗？难道你们不相信这些事是天主的讯息吗？这些是天主送来的书信；因此，让我们怀着相称的敬意进入圣堂，并怀着敬畏倾听这里所说的话。\par

\Quote{如果我不听人讲论，我进来做什么呢？}你说。这正是一切败坏和毁灭的根源。因为何需有人讲论呢？这需要出于我们的懈怠。因此，何需讲道呢？在\WorkTitle{圣经}中，一切事都是清楚明白的；必要的事都是显而易见的。但是因为你们是寻求快感的听众，所以你们也追求这些东西。请告诉我，\Person{保禄}讲话时用了什么华丽的辞藻？然而他却转变了世界。或者，不识字的\Person{伯多禄}用了什么辞藻？但你说，我不知道\WorkTitle{圣经}里所包含的内容。为什么呢？难道它们是用希伯来语讲说的吗？是用拉丁语，还是用外邦语言？难道不是用希腊语吗？但你说，它们表达得隐晦。什么是隐晦的？告诉我。难道没有历史吗？因为（当然）你们知道明显的部分，所以你们才询问隐晦的部分。\WorkTitle{圣经}里有无数历史。告诉我其中一件。但你们不能。这些是借口，只是空话。每天，你说，你听到同样的东西。那么告诉我，你在剧院里不也听到同样的东西吗？你在赛马场不也看到同样的东西吗？难道一切不都是相同的吗？升起的太阳不总是同一个吗？我们享用的食物不也是同样的吗？我愿意问你，既然你说每天听到同样的东西，那么告诉我，所读的段落是出自哪位先知？出自哪位宗徒，或哪封书信？但你却不能告诉我——你似乎是听到奇怪的东西。所以当你想要懈怠时，你说它们是同样的东西；但当你被质问时，你就像一个从未听过它们的人。如果它们是同样的，你就应该知道它们。但你却对它们一无所知。\par

这种情况值得哀叹——值得哀叹和抱怨：因为铸币者白白铸造了。正因为如此，你们更应当留意，因为它们是同样的东西，因为我们没有给你们带来辛劳，也没有讲论奇怪或变化的东西。那么，既然你们说那些是同样的东西，而我们的讲论却不同，我们总是对你们讲论新事，你们留意这些吗？绝不。但如果我说，为什么连这些你们也不记住呢？\Quote{我们只听过一次，}你们说，\Quote{怎么能记住呢？}如果我们说，为什么你们不留意那些其他的事呢？\Quote{总是说同样的东西，}你们说——而每一点都是懈怠和借口的言辞。但它们不会永远有用，因为终将有一个时候我们会徒然哀叹而毫无效果。愿天主禁止这事，并恩赐我们在此悔改，以理解力和对天主的敬畏来留心所讲的话，从而使我们被催促去适当地行善工，并以全备的勤勉修正我们自己的生活，好能获得那应许给爱祂的人们的福佑，借着恩宠和慈爱，等等。\par

\SourceRecord{Translated by John A. Broadus.；Nicene and Post-Nicene Fathers, First Series；Vol. 13.；Edited by Philip Schaff.；Buffalo, NY: Christian Literature Publishing Co.,；1889.；<http://www.newadvent.org/fathers/23053.htm>.}

% structure: p0001 source=h1 target=section reason=page-title

\section{《得撒洛尼后书》第四篇讲道}

\BibleRef{得后 2:6-9}\par

\Emph{\BibleQuote{现在你们也知道那阻止者，为要使他到了自己的时候才被显露。因为那不法之奥秘已经运行，只等到那现今阻止者被除去。那时，那不法者将显露出来，主耶稣要以口中的气息杀死他，并以自己来临的显现使他归于乌有；那不法者的来临，是依照撒殚的作为。}}\par

人自然会查问，那阻止者是什么，而后又想知道，为什么保禄表达得如此隐晦。那么，这阻止他显露的是什么？有些人说是圣神的恩宠，但另一些人说是罗马帝国，我大多数同意后者。为什么呢？因为如果他是指圣神，他就不会说得隐晦，而是明白地说：即现今圣神的恩宠，也就是神恩，阻止了他。而且，如果他要等待神恩停止时才来，那么他早就该来了，因为神恩早已停止了。但因为他说的是关于罗马帝国，他自然就暗指它，并且说得隐晦而模糊。因为他不想给自己招来多余的敌意和无用的危险。因为如果他说罗马帝国不久将瓦解，他们就会立刻把他当作瘟疫般的人，连同所有信徒，视为为此而生活和战斗的人，而压倒他们。而且他没有说它很快会来，尽管他总是这样说——但他说什么？\BibleQuote{为要使他到了自己的时候才被显露}，他说，\par

\BibleQuote{因为那不法的奥秘已经运行。}这里他提到尼禄，好像他是假基督的预像。因为尼禄也想要被当作神。他很好地说\BibleQuote{奥秘}，意思是不公开地运行，不像另一个那样，也不是毫无羞耻。因为如果在那个时代之前就有一个人，他说，在邪恶上不亚于假基督，那么现在若出现一个，又有什么奇怪呢？但他也不愿意清楚地指出他来：这不是出于怯懦，而是教导我们，在没有必要的情况下，不要给自己招来不必要的敌意。所以他在这里也这样说。\BibleQuote{只等到那现今阻止者被除去}，就是说，当罗马帝国被除去时，那时他就要来。这是自然的。因为只要对这个帝国的恐惧存在，就没有人愿意自高自大；但帝国瓦解后，他将攻击无政府状态，并试图夺取对人和天主的统治。因为正如之前的各个王国被消灭一样，例如，玛待王国被巴比伦人消灭，巴比伦人又被波斯人消灭，波斯人又被马其顿人消灭，马其顿人又被罗马人消灭；同样，这个帝国也将被假基督消灭，而假基督又被基督消灭，那时再没有阻止者了。这些事情，达尼尔非常清楚地启示了我们。\par

\BibleQuote{那时，}他说，\BibleQuote{那不法者将显露出来。}然后呢？安慰就在眼前。\BibleQuote{主耶稣要以口中的气息杀死他，并以自己来临的显现使他归于乌有；那不法者的来临，是依照撒殚的作为。}\par

因为火只是临近，甚至在到达之前，就使远处的微小动物麻痹并烧毁；同样，基督也仅凭他的命令和来临。他只需出现，这一切就被消灭。他仅凭显现，就遏止这欺骗。那么，这个依照撒殚作为而来的人是谁呢？\BibleQuote{具有一切能力，}但没有一样是真实的，都是为了欺骗。\BibleQuote{和虚假的奇事，}他说，即是虚假的，或引人走向虚假的。\par

第10节。\BibleQuote{并以一切不义的欺骗，为那丧亡的人。}\BibleRef{得后 2:10}\par

那么，你说，天主为什么允许这事？这是什么安排？若他的来临是为了毁灭我们人类，那有什么益处？不要害怕，亲爱的，且听他说：\Quote{在丧亡的人身上}，他才有能力，这些人即使他不来，也不会相信。那么有什么益处呢？就是这些丧亡的人本身将被封住口。如何封住？因为即使他来或不来，他们都不会相信基督；所以他来是为定他们的罪。为使他们没有机会说，因为基督说他是天主——虽然他从未公开这样说——但因为他以后的人宣告了这事，我们没有相信。因为我们听说只有一位天主，万物来自他，所以我们没有相信。假基督将除去他们的这个借口。因为当他来了，且命令的不是什么好事，而全是不法的事，却因虚假的徵兆而被人相信，这将封住他们的口。因为如果你们不相信基督，你们更不应该相信假基督。因为前者说他从圣父那里被派遣，后者则相反。为此基督说：\BibleQuote{我因我圣父之名而来，你们却不接待我；若别人因自己的名而来，你们反而接待他。}\BibleRef{若 5:43}但我们看见了徵兆，你们说。但在基督身上也行了许多大徵兆；你们更应相信他才对。然而关于这假基督，有许多预言指出他是那不法者，是丧亡之子，他的来临是依照撒殚的作为。但关于基督，预言的却是相反：他是救主，他带来无数的祝福。\par

第10至12节。\BibleQuote{因为他们没有接受对真理的爱，好使他们得救；为此，天主给他们派遣一种错误的效力，使他们相信谎话，使一切不信真理、反而喜爱不义的人都受审判。}\BibleRef{得后 2:10}\par

\Quote{好使他们受审判。}祂没有说，好使他们受惩罚；因为在这之前，他们已要受惩罚；而是说\Quote{好使他们被定罪}，也就是说，在那可畏的审判座前，使他们无可推诿。\Quote{那些不信真理，反喜爱不义的人。}祂称基督为\Quote{真理之爱}。因为经上说：\Quote{因为他们没有接受真理之爱。}原来祂既是爱，又是真理，并为两者而来，既出于爱人，又为真道之事。\par

\Quote{反而喜爱}，祂说，\Quote{不义。}因为他来是为毁灭人，伤害他们。那时他还有什么事情做不出来呢？他要藉他的命令和人们对他的畏惧，改变并扰乱一切。他将因他的权能、他的残暴、他的不法命令而处处可畏。\par

但是不要害怕。\Quote{在那些灭亡的人身上}，他将显出自己的能力。因为那时厄里亚也要来，使信者安心，这也是基督所说的：\BibleQuote{厄里亚要来，且要重整一切。}\BibleRef{玛 17:11}因此经上说：\BibleQuote{带着厄里亚的精神和能力。}\BibleRef{路 1:17}因为他没有像厄里亚那样行神迹奇事。因为经上说：\Quote{若翰没有行过奇迹，但若翰关于这人所说的一切都是真的。}那么，\Quote{带着厄里亚的精神和能力}这是怎样呢？就是说，他要承担同样的职务。正如前者是祂第一次来临的先行者，同样，后者将是祂第二次荣耀来临的先行者，为此他已被保留。所以我们不要害怕。祂已平静了听众的心。祂使他们不再认为现在的事可怕，而是值得感恩。因此他又说：\par

第13节。\BibleQuote{但我们常该为你们感谢天主，主所爱的弟兄们，因为天主从起初拣选了你们，借着圣神的圣化及对真理的信仰，而得救恩。}\BibleRef{得后 2:13}\par

如何得救？是藉着圣神圣化你们。因为这些是我们得救的有效原因。绝不是出于行为，绝不是出于义行，而是出于对真理的信仰。这里再次用\Quote{在……内}表示\Quote{藉着}。他说：\Quote{并藉着圣神的圣化，}\par

第14节。\BibleQuote{为此祂藉着我们的福音召叫了你们，为获得我们的主耶稣基督的光荣。}\BibleRef{得后 2:14}\par

这也不是小事，如果基督将我们的得救视为祂的光荣。因为爱人的朋友的光荣，就在于得救的人众多。那么我们的主该是何等伟大，如果圣神如此渴望我们的得救。为什么他不先说信仰？因为即使在圣化之后，我们还需要大量的信德，以免动摇。你看到祂如何表明一切都非出于他们自己，而是出于天主吗？\par

第15节。\BibleQuote{所以，弟兄们，你们要站立稳定，持守你们所学习的传统，无论是藉我们的话语，或是藉我们的书信。}\BibleRef{得后 2:15}\par

由此可见，他们并非全部通过书信传授，还有许多未成文的，而且两者同样值得信赖。所以我们也当认为教会的传统是值得信赖的。这是一个传统，不必再寻求其他。这里他表明有许多人动摇了。\par

第16-17节。\Quote{愿我们的主耶稣基督自己，以及我们的天主父，祂爱了我们，并藉恩宠赐予我们永久的安慰和美好的希望，安慰你们的心，并使你们在一切善言善行上坚定稳固。}\par

劝诫之后又有一番祈祷。因为这才真正有益。他说：\Quote{祂爱了我们，并藉恩宠赐予我们永久的安慰和美好的希望。}如今那些贬低圣子的人在哪里？因为祂在圣洗的恩宠中被列在圣父之后？看，这里恰恰相反。他说：\Quote{祂爱了我们，并赐予我们永久的安慰。}这安慰是怎样的呢？就是未来事物的希望。你看到吗？祂藉着祈祷的方式激励他们的心灵，将天主不可言喻的眷顾作为抵押和标记赐给他们。他说：\Quote{安慰你们的心，在一切善言善行上，}也就是说，藉着一切善言善行。因为基督徒的安慰就在于做中悦天主的好事。看，祂如何降卑他们的心。他说：\Quote{祂赐予我们安慰和美好的希望，藉着恩宠。}同时，祂也使他们充满对未来事物的美好希望。因为如果祂藉恩宠赐予了这么多，何况将来的事呢？祂说：我固然说了，但一切都是出于天主。\Quote{坚定}，即坚固你们，使你们不被动摇，也不转向。因为这是祂的工作，也是我们的工作，因此它既在教义上，也在行动上。因为安慰就是被坚定。因为当一个人不转向时，他就能以极大的耐心承受一切所遇；但如果他的心意动摇，他就再也不能做出任何善行或高尚的行动，而像一个双手瘫痪的人，他的灵魂也摇动，因为他不完全确信自己正迈向某个美好的终点。\par

第三章第1节。\BibleQuote{此外，弟兄们，请为我们祈祷，好使主的圣言得以迅速传开，并受光荣，如同在你们中间一样。}\BibleRef{得后 3:1}\par

他确实曾为他们祈祷，使他们得以坚定；现在他请求他们，恳求他们为他祈祷，不是为使他免遭危险——因为他正是为此被派——而是为\Quote{使主的圣言得以迅速传开，并受光荣，如同在你们中间一样。}这一请求还伴随着称赞。\Quote{如同在你们中间一样。}\par

第2节。\BibleQuote{并使我们能摆脱无理和邪恶的人，因为不是人人都有信德。}\BibleRef{得后 3:2}\par

这是那人显示他自己的危险的方式，他尤其为此恳求他们。\Quote{从无理和邪恶的人}，他说，\Quote{因为众人没有信德}。因此他是在谈论那些反对宣讲的人，他们反对并争辩教义。为此他通过说\Quote{因为众人没有信德}来暗示。在我看来，他这里并未着眼于危险，而是着眼于那些反驳和阻碍他言语的人，如同铜匠\Person{亚历山大}一样。因为他说，\BibleQuote{他极力抵挡我们的话}\BibleRef{弟后 4:15}。也就是说，有些人是被赋予的。好像他在谈论父亲的遗产，说\Quote{不是所有人都能在王宫中服务}。同时，他也激励他们，因为他们已经拥有这样的信心基础，既能拯救他们的老师脱离危险，又能促进他的宣讲。\par

因此，我们也说同样的事。不要让任何人因骄傲而谴责我们，也不要因过度的谦卑而剥夺我们如此巨大的援助。因为我们说话并非出于\Person{保禄}所说的同一动机。他确实说这些事是为了安慰他的门徒；而我们是为了收获某些巨大而美好的果实。我们非常确信，如果你们都愿意同心合意地为我们的卑微向天主伸出双手，你们将在一切事上成功。因此，让我们用祈祷和恳求与我们的敌人作战。因为如果古代的人用武器与人作战，那么我们就更应当这样与手无寸铁的人作战。\Person{希则克雅}就这样战胜了\Place{亚述}王，\Person{梅瑟}战胜了\Place{阿玛肋克}，\Person{撒慕尔}战胜了\Place{阿市刻隆}人，\Place{以色列}战胜了三十二个王。如果在需要武器、战阵和战斗的地方，他们放下武器，转而祈祷；那么在这里，事情必须单靠祈祷来完成，我们岂不更应该祈祷吗？\par

但在那里，你说，统治者为人民求情，但你请求人民为统治者求情。我也知道。因为那时被统治的人是可怜和卑微的。因此，他们仅凭着他们统帅的功绩和德行而得救；但现在，当天主的恩宠得胜，我们将在被统治的人中发现许多人，或者更确切地说，大部分人在很大程度上超越了他们的统治者时，不要剥夺我们的这种援助，举起我们的手，使它们不倦怠，为我们开口，使口不闭合。恳求天主——为此恳求祂。这确实是为我们而做的，但完全是为你们。因为我们被任命为你们的益处，我们关心你们的利益。你们每个人，私下和公开地，都恳求吧。看\Person{保禄}说：\BibleQuote{因了许多人的代求，而赐予我们的恩赐，藉着许多人的感谢，也为我们而归于天主}\BibleRef{格后 1:11}；也就是说，使祂赐恩宠给许多人。如果在人的情况下，人民上前请求赦免那些被定罪并被带去处决的人，而国王因着群众的缘故撤销判决，那么天主更会因着你们——不是因你们的数量，而是因你们的德行——而受到影响。\par

因为我们所拥有的敌人是凶猛的。你们每个人都焦虑地考虑自己的利益，但我们考虑所有人的利益。我们站在战斗中被压迫的部分。魔鬼更猛烈地武装起来攻击我们。因为在战争中，对面的人首先试图推翻将军。为此，他所有的战友都急忙赶到那里。为此，有巨大的骚动，每个人都努力拯救他；他们用盾牌围绕他，希望保护他的身体。听所有人民对\Person{达味}说的话（我说这些，并非将自己比作\Person{达味}，我没有那么疯狂，而是因为我希望显示人民对他们的统治者的爱戴）。他们说：\BibleQuote{你不可再与我们一同出战，免得你熄灭以色列的明灯}\BibleRef{撒下 21:17}。看他们多么渴望保护这位老人。我极其需要你们的祈祷。正如我所说，不要让任何人因过度的谦卑而剥夺我的这种联盟和援助。如果我们的部分获得嘉许，你们自己的部分也会更荣耀。如果我们的教导丰盛地流淌，财富将归于你们。听先知说：\BibleQuote{牧人岂是牧养自己呢？}\BibleRef{则 34:2}\par

你看到\Person{保禄}不断寻求这些祈祷吗？你听到\Person{伯多禄}是如何从监狱中被释放的，当时有人为他热切祈祷？\BibleRef{宗 12:5} 我确信，你们的祈祷，如此同心合意地献上，将产生巨大的效果。你们不认为，让我的卑微接近天主并为他如此众多的人民恳求，是太过重大的事吗？因为如果我没有为自己祈祷的信心，就更不能为他人祈祷了。因为恳求天主怜悯他人，属于那些有崇高评价的人；属于那些使天主对自己慈祥的人。但一个犯罪者本人，怎能替别人恳求呢？但尽管如此，因为我以父亲的心肠拥抱你们，因为爱敢于一切，不仅在教会内，也在家中，我首先为你们灵魂和身体的健康祈祷。因为在肉体上的孩子面前，没有其他人民。如果\Person{约伯}立刻为他在肉身上的儿女献上如此多的祭献，那么我们更应当为我们的属神子女这样做。\par

为什么我说这些事？因为如果我们这些离这伟大工作如此遥远的人，尚且为你们呈上恳求和祈祷，那么你们更应当这样做。因为一个人为许多人恳求，是极其大胆的，需要很大的信心；但许多人聚集在一起为一个人恳求，则毫无负担。因为每个人这样做，不是依靠自己的德行，而是依靠众多的人及其同心合意，天主处处对此十分尊重。因为祂说：\BibleQuote{无论在哪里，有两三个人因我的名聚在一起，我就在他们中间。} \BibleRef{玛 20:18} 如果两三个人聚集在一起，祂就在他们中间，那么祂更在你们中间。因为一个人独自祈祷所不能得到的，他同众人一起祈祷就会得到。为什么呢？因为虽然他自己的德行没有，但共同一致的力量很大。\par

经上说：\BibleQuote{无论在哪里，有两三个人，} …… \BibleQuote{聚集在一起。} 为什么说\Quote{两}呢？因为如果有一个人因你的名，为什么你不在那里呢？因为我愿意大家都聚在一起，不要分离。因此，让我们紧密相连；让我们在爱中彼此结合，不要让任何人分开我们。如果有人指责或受到冒犯，不要让他记在心里，无论是对邻人还是对我们。我请求你们这样做：到我们这里来，提出指控，听取我们的辩护。经上说：\BibleQuote{你要责备他，} 经上说，\BibleQuote{恐怕他没有说这话。你要责备他，恐怕他没有做这事；} \BibleRef{德 19:14} 如果他做了，使他不再加添。因为我们或者为自己辩护，或者被定罪后请求宽恕，并从此努力不再犯同样的过错。这对你们和我们都有益处。因为你们可能无缘无故地指责我们，但当你们知道事情的真相后，就会纠正自己；而我们无意中冒犯了你们，也会得到纠正。因为对你们来说，这并不有益。因为惩罚是为那些说闲话的人预备的。但我们放下指控，无论是假的还是真的。假的，通过证明它们是假的；真的，通过不再做同样的事。因为管理这么多事的人难免会有疏忽，并因无知而犯错。如果你们每个人有一所房子，管理妻子、儿女、奴仆，有的多有的少，在如此容易数算的灵魂中，尚且不得不非自愿地犯许多错误，或因无知，或因想要纠正某事；那么管理这么多人的我们更是如此。\par

愿天主仍使你们增多，祝福你们，小的与大的并存！因为虽然随着人数的增加，关怀变得更重，但我们仍不停祈祷，愿我们的关怀得以增加，愿这数目得以增添，增多许多倍，没有穷尽。因为父亲们虽然常被子女的数目所烦扰，却不愿失去任何一个。我们和你们之间一切平等，甚至我们祝福中最首要的也是如此。我领受圣体圣事并不比你们更多，你们也不更少，而是同样分享同一圣体。即使我先领受，也不是什么特权，因为即使在孩童中，年长的先伸手赴宴，但并无任何优势。但我们中间一切平等。那滋养我们灵魂的救恩之生命，同样荣耀地赐予我们双方。我并非领受一只羔羊，你们领受另一只，而是我们领受同一只。我们都有同一圣洗圣事。我们被赐予同一圣神。我们都奔向同一天国。我们都是基督的兄弟姐妹，我们共同拥有一切。\par

那么我的优势在哪里？在挂虑、劳苦、焦虑、为你们忧伤。但没有什么比这忧伤更甜美，因为即使母亲为她的孩子忧伤，她也因忧伤而喜悦，她仔细惦记自己所生的孩子，她因自己的挂虑而喜悦。然而挂虑本身是苦涩的，但若为了孩子，至少其中有许多快乐。我生了你们许多人，但之后是我的痛苦。就血肉的母亲而言，痛苦在先，然后才是生产。但在这里，痛苦持续到最后一息，免得在生产之后还有堕胎。而我们有更深切的渴望；因为虽然或许别人生了你们，但我仍被挂虑所困。因为我们不是自己生你们，而是全赖天主的恩宠。但如果我们都藉着圣神而生，那么称我所生的人为他的孩子，称他所生的人为我的孩子，这并没有错。所以你们要思考这一切，并伸出你们的手，使你们成为我们的光荣，我们也成为你们的光荣，在我们主耶稣基督的日子，愿天主恩赐我们都能怀着信心见到那一天，藉着我们的主耶稣基督，愿光荣归于祂，等等。\par

\SourceRecord{Translated by John A. Broadus.；Nicene and Post-Nicene Fathers, First Series；Vol. 13.；Edited by Philip Schaff.；Buffalo, NY: Christian Literature Publishing Co.,；1889.；<http://www.newadvent.org/fathers/23054.htm>.}

% structure: p0001 source=h1 target=section reason=page-title

\section{得撒洛尼后书讲道第五篇}

得撒洛尼后书 3:3-5\par

\Emph{\Quote{但是主是信实的，祂必坚固你们，保护你们脱离那恶者。我们靠着主深信你们现在和将来都会遵行我们所命令的事。愿主引导你们的心进入天主的爱和基督的忍耐。}}\par

我们既将一切托付于圣人们的祈祷，也不应自己闲散，陷于恶行，无所作为；也不应在行善时轻视那助佑。因为为我们所作的祈祷的确能成就大事，但当我们自己也努力时。为此，\Person{保禄}也为他们祈祷，又藉恩许给他们保证，说：\Quote{但是主是信实的，祂必坚固你们，保护你们脱离那恶者。}因为如果祂选择了你们得救恩，祂就不欺骗你们，也不容许你们完全丧亡。但为了不使他们因此而懈怠，并免使他们以为全赖天主而自己睡觉，请看他是如何也要求他们合作，说：\Quote{我们靠着主深信你们现在和将来都会遵行我们所命令的事。}他说，\Quote{主}确实是\Quote{信实的}，既应许了拯救，就必拯救；但正如祂所应许的。祂是如何应许的呢？只要我们愿意，并听从祂；不单是听从，也不像木头石头一样，不动。\par

他巧妙地引入了这话：\Quote{我们靠着主深信}，就是说，我们信赖祂的慈爱。他又使他们谦卑，使一切都依赖于此。因为如果他说：我们深信你们，那称赞固然很大，但却不会教导他们万事依赖天主。如果他说：我们信赖主，祂必保护你们，而没有加上\Quote{至于你们}和\Quote{你们现在和将来都会遵行我们所命令的事}，他就会使他们更加懈怠，将一切推给天主的能力。因为我们固然应将一切托付于祂，但也要自己努力，投身于劳作和争战。他表明，即使我们的德行足以拯救我们，它仍需持之以恒，与我们同在，直到我们最后一息。\par

\Quote{但是主，}他说，\Quote{引导你们的心进入天主的爱和基督的忍耐。}\par

他又称赞他们，并祈祷，表示对他的关心。因为当他即将开始责备时，他先安抚他们的心，说：\Quote{我相信你们会听从}，并请求他们的代祷，又再次向他们祈求无限的祝福。\par

\Quote{但是主}，他说，\Quote{引导你们的心进入天主的爱。}因为有许多事使我们偏离爱，有许多道路将我们从那里引开。首先是\OriginalTerm{玛门}{Mammon}的道路，它仿佛将某种无耻的手放在我们的灵魂上，紧紧抓住它，把它拖走，甚至违背我们的意志。其次是虚荣，以及常常的苦难和诱惑，使我们偏离。为此，我们需要天主的助佑，如同某种风，推动我们的帆，如同被强风，进入天主的爱。不要对我说：\Quote{我爱祂，甚至超过爱我自己。} 这是空话。用你的行为向我证明，如果你爱祂超过爱你自己。为了天主而轻看钱财，然后我就相信你爱祂甚至超过爱你自己。但是，你为了天主而不轻视财富，怎么能轻视你自己呢？为何我说财富？你连贪恋都不能轻视——这本是即使没有天主的诫命你也应当做的——又如何能轻视你自己呢？\par

\Quote{和基督的忍耐}，他说。\Quote{忍耐}是什么意思？是叫我们如同祂一样忍耐，或是我们行那些事，或是我们也要用忍耐等候祂，就是说，我们应当准备好。因为祂既应许了许多事，并且祂亲自来审判活人死人，让我们等候祂，让我们忍耐。但无论他在何处提到忍耐，当然意味着苦难。因为这就是爱天主：忍受，而不烦恼。\par

第6节。\Quote{现在我们奉主耶稣基督的名命令你们，弟兄们，要远离一切行事不守规矩、不遵从我们所传给你们的传统的弟兄。}\BibleRef{得后 3:6}\par

这就是说，不是我们说的这些话，而是基督，因为\Quote{奉主耶稣基督的名}的意思就是\Quote{藉着基督}。显示信息的可怕性，他说，藉着基督。因此基督命令我们绝不可闲散。\Quote{要远离}，他说，\Quote{每一个弟兄。}不要对我提起富人，不要对我提起穷人，不要对我提起圣人。这是不规矩。\Quote{行事}，他说，即生活。\Quote{不遵从我们所传给你们的传统。}传统，他说，是通过行为的传统。他常恰当地称之为传统。\par

第7-8节。\Quote{你们自己知道应当怎样效法我们：因为我们在你们中间并没有不守规矩，也没有白吃任何人的饭。}\par

即使他们吃了，也不是白吃。\Quote{因为工人}，他说，\Quote{配得他的工价。}\BibleRef{路 10:7}\par

\Quote{但我们在劳苦和辛劳中，昼夜工作，免得你们中任何人受累。并不是因为我们没有权利，而是要以自己给你们作榜样，叫你们效法我们。因为我们从前在你们那里的时候，曾吩咐你们说，若有人不肯工作，就不可吃饭。}\par

请看他在前书论及这些事时，语气确实较为温和。例如，他说：\BibleQuote{弟兄们，我们劝你们……你们更要不断进步……并要用心}——见：\BibleRef{得前 4:1}——他绝未说：\Quote{我们命令}，也未说\Quote{因我们的主耶稣基督的名}，这语气令人畏惧且暗含危险；而是说\BibleQuote{你们更要进步}，他说，\BibleQuote{并要用心}，这是一个劝勉人修德者的话；他说：\BibleQuote{使你们行事端庄（得体）}——见：\BibleRef{得前 4:12}。但这里却无此类语气，他说：\BibleQuote{若有人不愿工作，就不可吃饭}。因为保禄并非出于不得已，他本有权利闲散，又承担如此伟大的工作，却仍然工作，且不仅工作，而是\Quote{昼夜劳碌}，以致还能帮助他人——何况其他人更应如此。\par

第11节。\BibleQuote{因为我们听说，你们中有些人行事不规矩，什么工也不作，反倒好管闲事。}——见：\BibleRef{得后 3:11}\par

他在这里确是这样说的；但在前书中，他说：\BibleQuote{使你们行事端庄，对待外人}——见：\BibleRef{得前 4:12}。出于什么原因？或许那时还没有这类事。因为在另一场合劝告时，他说：\BibleQuote{施比受更为有福}——见：\BibleRef{宗 20:35}。但\Quote{行事端庄}这句话并不涉及不守规矩；因此他加上了：\BibleQuote{使你们不至于缺少什么}——见：\BibleRef{得前 4:12}。在这里他又提出了另一个必要性，即如此行是光荣而美好的，对众人都是有益的。（因为他在继续讲论时说：\BibleQuote{行善不可厌倦}。）因为那些闲散而又有能力工作的人，必定会好管闲事。但施舍只应给予那些不能靠双手劳动养活自己的人，或者那些从事教导、全心致力于教导工作的人。他说：\BibleQuote{不可笼住牛的嘴，当它踹谷的时候}——见：\BibleRef{申 25:4}，\BibleQuote{工人得工资是应当的}——见：\BibleRef{弟前 5:18}。因此，他也不闲散，而是接受工作的报酬，且是极大的工作报酬。但闲散地祈祷和斋戒，并非双手的工作。因为他在这里所说的工作是双手的工作。为了使你不致怀疑任何这类事，他加上了：\par

\BibleQuote{什么工也不作，反倒好管闲事。我们因主耶稣基督吩咐并劝勉这样的人。}\par

因为他严厉地触动了他们，为了使言论更温和，他加上了\BibleQuote{因主}，这又是权威而令人敬畏的。\par

他说：\BibleQuote{要安静，做工，吃自己的饭。}\par

为什么他没有说：如果他们不守规矩，就由你们供养；而是要求两样：他们既要安静，又要做工？他说：\Quote{为使他们吃自己的饭}，而不是别人的饭。\par

第13节。\BibleQuote{但是你们，弟兄们，行善不可厌倦。}——见：\BibleRef{得后 3:13}\par

请看慈父的心肠如何立即被触动。他不能再继续谴责下去，而是再次怜悯他们。且看他的智慧！他没有说：\Quote{宽恕他们，直到他们改正}；而是说什么？\BibleQuote{但是你们，行善不可厌倦。}他说：你们远离他们，并责备他们；但不可让他们因饥饿而死。那么，他说，如果他从我们这里得到充裕的物资，却仍然闲散呢？在这种情况下，他说，我已谈到一种温和的补救办法：你们远离他，即不要与他自由交往；表明你们是愤怒的。这不是小事。因为这样的责备给予一个弟兄，如果我们真的想改正他，是非常有效的。我们并非不知道责备的方法。请告诉我，如果你有一个亲弟兄，你会眼睁睁看着他因饥饿而消瘦吗？我真的不这样认为；也许你甚至会纠正他。\par

第14节。\BibleQuote{如果有人不听从我们在这封信中的话。}——见：\BibleRef{得后 3:14}\par

他没有说：\Quote{那不听的人，是不听我}，而是说\BibleQuote{要记下那人}。这不是轻微的惩罚。\BibleQuote{不要与他交往。}然后他又说：\BibleQuote{使他羞愧。}并且他不允许他们走得更远。因为正如他所说：\BibleQuote{若有人不愿工作，就不可吃饭}，因怕他们饿死，他又加了：\BibleQuote{行善不可厌倦。}这样，他说：\Quote{你们要远离，不要与他交往}，然后又怕这行为本身会使他脱离弟兄情谊——因为自暴自弃的人若不被允许自由交谈，很快就会迷失——他加上了：\par

第15节。\BibleQuote{但不要以他为仇人，而要劝他如弟兄。}——见：\BibleRef{得后 3:15}\par

由此他表明，他给予那人沉重的惩罚，即剥夺他自由交谈的权利。\par

因为如果与许多人一同接受施舍尚且可耻，那么当他们一边施舍一边斥责、并且回避时，这责难何其大，足以刺痛灵魂。因为如果只是迟延给予、并且发怨言，他们便激怒了受施者——请不要提那些厚颜的乞丐，而是指信徒——如果他们边给予边斥责，还有什么做不出来？这岂不是与某样惩罚相等？我们却不这样做，反而如同受了极大伤害，对求我们的人百般侮辱和拒绝。你不给也就罢了，为什么还要使他悲伤？他说：\Quote{要劝诫他们，如同弟兄}，不要像对待仇敌一样侮辱他们。劝诫自己弟兄的人不会公开去做，他不会张扬侮辱，而是私下并非常巧妙地去作，如同受伤、哭泣哀叹一般。因此让我们以弟兄的心态施与，以弟兄的善意劝诫，不是因为给予而痛苦，而是因为他违反诫命而痛苦。因为有什么好处呢？即使给了之后又侮辱，你摧毁了施与的喜乐。但当你既不给又侮辱时，你对那可怜不幸的人做了什么亏待呢？他到你这里来，期望从你得到怜悯，却带着致命的打击离去，哭得更厉害了。因为当他因贫穷被迫行乞，又因行乞而受侮辱时，你想那些侮辱他的人将受何等大的惩罚？圣经说：\BibleQuote{欺压穷人的，辱没造他的主。}\BibleRef{箴 14:12}。请问，祂使他为你成为穷人，好使你能够治愈自己——你却侮辱那为你成为穷人的人？这是何等顽固！何等忘恩负义！他说：\Quote{要把他当作弟兄劝诫}，并且在你给予之后，他命令你劝诫他。但如果连不给予时也要侮辱他，我们还有什么借口？\par

\Quote{愿平安的主亲自在各时各处赐给你们平安。}\BibleRef{得后 3:16}\par

请看，当他提到当做的事时，他用祈祷来标示它们，加上祈求与恳祷，如同在某些珍藏品上加上标记。他说：\Quote{赐给你们平安，在各时各处}。因为很可能从这些事情中会产生争执——那些人被激怒，而其他人不像从前那样乐意供给他们——所以他很有理由现在为他们献上这祈祷，说：\Quote{在各时}赐给你们平安。这正是所祈求的，好使他们永远拥有它。他说：\Quote{在一切方式上}。\Quote{在一切}是什么意思？就是使他们从任何方面都没有争斗的缘由。因为平安处处都是好事，甚至对那外人也是如此。听他在别处说：\BibleQuote{若是可能，尽力与众人和睦。}\BibleRef{罗 12:18}。因为没有什么比和平、不受干扰、远离所有仇恨、没有敌人更有利于正确完成我们所希望的事了。\par

\Quote{愿主与你们众人同在。}\par

\Quote{我\Person{保禄}亲笔问候，这是我在每封信上所用的记号：我这样写。愿我们的主耶稣基督的恩宠与你们众人同在。}\BibleRef{得后 3:17-18}\par

他说他在每一封信上都这样写，好使没有人能伪造它们，因为附有他的署名作为重要的记号。他称这祈祷为问候，表明他们那时所做的一切都是属灵的；即使是在应当问候的时候，此事也伴随着益处；这是祈祷，不仅是友谊的象征。他以此开始，也以此结束，用坚固的城墙保卫了他处所说的话，奠定了安全的基础，也带它到达安全的终点。他说：\Quote{愿恩宠与平安归于你们}；又说：\Quote{愿我们主耶稣基督的恩宠与你们众人同在。阿们。}这也就是主所应许的，祂对门徒说：\BibleQuote{看，我天天与你们在一起，直到世界的终结。}\BibleRef{玛 28:20}。但这事发生在我们愿意的时候。因为如果我们自己远离，祂不会完全与我们同在。祂说：\Quote{我必与你们同在，直到永远。}所以我们不要赶走恩宠。他吩咐我们远离每一个行事素乱无章的弟兄。这在当时是极大的恶事——被隔绝于整个弟兄团体之外。他以此惩罚所有人，如同在致格林多人书里所说的：\BibleQuote{不可与他一同吃饭。}\BibleRef{格前 5:11}。但现在大多数人并不认为这是大恶。但一切都被混淆和败坏了。我们随意地与奸夫、淫乱者、贪婪之人交往，视为理所当然。如果我们应当远离只靠懒惰养活的人，何况远离其他人呢？并且，为使你知道被隔绝于弟兄团体是何等可怕的事，以及这给以正确心态接受责罚的人带来何等益处，请听那位因罪恶而骄傲、行至邪恶极致、行连外邦人中也不曾提到的淫乱之人，并且对自己的创伤毫无知觉——因为这正是极端悖逆——他后来却被折服和谦卑，以至于\Person{保禄}说：\BibleQuote{这样的人，受了多数人的责罚，已经够了；所以你们应当向他重申你们的爱。}\BibleRef{格后 2:6}。因为那时他就像一个与身体分离的肢体。\par

但原因，以及当时使这件事如此可怕的原因，是因为甚至与他们同在也被他们视为极大的祝福。因为如同人们同住一屋、同属一父、同享一桌，那时他们也是这样居住在每一个教会中。因此，从如此大的爱中堕落是何等大的祸患！但现在，它甚至不被认为是大祸患，因为当我们彼此联合时，也根本不被视为大事。那时属于惩罚次序的事，如今因爱之冷淡，甚至在惩罚之外就发生了，我们无缘无故地、且因冷淡而彼此疏远。因为没有爱，就是一切祸患的原因。这爱已解开了所有纽带，丑化了教会中一切可敬且荣耀的事，我们本应以这些为荣。\par

导师的信心是大的，当他因自己的善行而有权责备门徒时。因此保禄也说：\BibleQuote{因为你们自己知道该怎样效法我们。}（\BibleRef{得后 3:7}）他应当是一位更以生活、而非以言语为师的导师。不可有人认为这话是出于夸耀之心。因为他是不得已才说的，且为了众人的益处。他说：\BibleQuote{因为我们没有在你们中间不守规矩，也没有白吃任何人的饭。}你难道看不出他的谦卑吗？因他把这称为\Quote{白吃}和\Quote{不守规矩的行为}？他说：\BibleQuote{我们没有在你们中间不守规矩，也没有白吃任何人的饭。}这里他显示他们或许也是贫穷的；并且你不要说他们是贫穷的。因为他所讨论的是穷人，以及那些除了自己双手的劳动之外别无生计的人。因为他并没有说，他们可以从他们的父亲那里获得，而是说他们应当亲手工作，吃自己的饭。因为他说，我这个教义之言的宣讲者，尚且害怕成为你们的负担，何况那不为你们服务的人呢？因为这确实是一个负担。当一个人不很乐意给予时，这也是一个负担；但这不是他所暗示的，而是好像他们不容易做到一样。你为什么不做工呢？因为天主赐给你双手，正是为此，不是要使你从别人那里接受，而是要使你给予别人。\par

但是他说：\Quote{愿主与你们同在。}如果我们做主的事，我们也可以为自己献上这祈祷。因为请听基督对祂的门徒说：\BibleQuote{你们要去使万民成为门徒，因父及子及圣神之名给他们授洗，教训他们遵守我所吩咐你们的一切。看！我与你们天天在一起，直到今世的终结。}（\BibleRef{玛 28:19}）如果你们做这些事，无疑地。因为那许诺不仅是对他们，也是对那些跟随他们脚步的人说的，这从祂所说的\Quote{直到今世的终结}是显然的。\par

那么主对不是教师的人说什么呢？你们每个人，若愿意，都是教师，虽然不能教导别人，却能教导自己。先教导你自己。如果你教训人遵守祂所吩咐的一切事，即使藉此方法，也会有许多人效法你。因为如盏灯，当它发光时，能点燃万盏灯，但若熄灭，连自己也不发光，也不能点亮其他灯；同样，在纯洁的生活上，如果我们里面的光在照耀，我们将使数不清的门徒和教师，因为我们被摆在面前作为他们效法的榜样。因为从我口中发出的言语，不能像你们的生活那样使听者受益。因为请告诉我，一个被天主所爱、在美德中发光、且有妻子的人（因为一个有妻子、儿女、仆人和朋友的人仍能取悦天主），难道他不能比我更使众人受益吗？因为他们每月会听我讲一次或两次，甚至一次也没有，即使他们将所听到的保留到教堂的门槛，他们很快就丢掉了；但他们不断看到那人的生活，就大得裨益。因为当受人侮辱时，他不还口，这岂不几乎将温顺的敬意灌输并铭刻在侮辱者的灵魂中吗？虽然他因愤怒而羞惭或困惑，没有立即承认受益，但无论如何他立即意识到了。一个傲慢的人，即使是一头野兽，与一个以善报恶的人交往，不可能不大大受益。因为虽然我们不行善，但我们却都称赞并钦佩善。再者，妻子若看见丈夫温和，常与他在一起，就大得裨益，孩子也是如此。因此，每个人都有能力成为教师。因为他说：\BibleQuote{你们要彼此扶持，正如你们现在所做的。}（\BibleRef{得前 5:11}）因为请告诉我，家庭是否遭受了什么损失？妻子因为更软弱、更奢侈、喜爱装饰而困扰；丈夫若是一位哲人，视损失如无物，既安慰她，又劝她刚强地承受。那么，请告诉我，他岂不是比我们的言语更使她受益吗？因为说话是容易的，但到了必须行动时，却处处困难。因此，人性往往更受行为的规范。而美德如此优越，以至于即使一个奴隶也常常与主人一同使整个家庭受益。\par

因为保禄并非徒然，也不是无理由地常常命令他们实践德行，并服从主人，与其说是注重主人的服务，不如说是为了使天主的话和道理不受亵渎。但当它不被亵渎时，很快也会受赞美。我知道许多家庭因奴隶的美德而大得裨益。但如果一个受权威管辖的仆人能够改善他的主人，那么主人更能改善他的仆人。那么，我恳求你们，与我分担这职务。我一般向所有人讲话，你们每个人私下里，并让每个人为自己邻人的得救负责。因为关于这些事情，一个人应当管理他的家人，请听保禄在哪里派遣妇女去受训导：\Quote{如他们愿意学习什么，让她们在家里问自己的丈夫} \BibleRef{格前 14:35}；他没有带领她们到导师那里。因为在学习的学校里，连门徒中也有教师，在教会中也是如此。因为他希望导师不被所有人烦扰。为什么？因为那样会有很大的好处，不仅导师的劳苦减轻，而且每个门徒经过努力，也能很快成为导师，专注于这件事。\par

请看妻子提供了多大的服务。她管理房屋，照管家中一切事务，统领她的婢女，亲手为她们制衣，使你被称为儿女的父亲，她使你远离妓院，帮助你贞洁生活，制止了自然的强烈欲望。你也应当回报她。如何做？在属灵的事情上伸出你的手。你听到的任何有益的事，像燕子一样，衔在嘴里，带到母亲和孩子口中。因为如果在其他事情上，你认为自己配得优先地位，占据头的地位，但在教导上却放弃你的位置，这岂不荒谬吗？统治者不应在荣誉上胜过被统治的人，而应在德行上。因为这正是统治者的职责，因为另一个是被统治者的部分，但这是统治者自己的成就。如果你享有许多荣誉，这对你无关紧要，因为你从别人那里得到它。如果你在德行上闪耀，这完全是你自己的。\par

你是女人的头，那么让头调节身体的其他部分。难道你没有看到，头不仅在位置高于身体，而且在远见上，像舵手一样引导整个身体？因为在头里有身体和灵魂的眼睛。由此流出看见的能力和指导的力量。身体的其他部分被指定为服务，而头被设定为发令。所有感官都从那里产生并有其源头。从那里发出言语的器官、视觉、嗅觉和所有触觉。因为神经和骨骼的根源也来自那里。你没有看到它在远见上优越于在荣誉上吗？让我们这样统治女人；让我们超越她们，不是通过寻求从她们那里得到更大的荣誉，而是通过她们从我们这里得到更大的益处。\par

我已经表明她们带给我们不小的益处，但如果我们愿意在属灵的事情上报答她们，我们就超越她们。因为在身体的事情上不可能提供等价物。为什么？你贡献大量财富吗？但却是她保存了财富，而她的这份关心是一种等价物，因此需要她，因为许多人拥有大量财产，却因为没有人照管而全部失去。至于孩子，你们共同参与，从每个人得到的益处是相等的。她反而在这些事情上有更辛苦的服务，总是孕育后代，并受生产之痛；所以只有在属灵的事情上你才能超越她。\par

因此，我们不要关心如何拥有财富，而要关心如何将托付给我们的灵魂充满信心地呈献给天主。因为通过规范他们，我们也最有益于自己。因为教导别人的人，即使他什么也不做，但在说话时也会感到愧疚，当他看到自己要对那些他责备别人的事情负责时。既然我们既有益于自己，也有益于他们，并通过他们有益于家庭，而这最中悦天主；让我们不要厌倦照顾我们自己的灵魂和那些服侍我们的人，以便我们为一切获得赏报，并带着许多财富到达我们的母亲圣城，那上面的耶路撒冷，愿天主使我们永不从那里坠落，反而在最卓越的生活道路上闪耀，得以充满信心地见到我们的主耶稣基督；愿光荣、权能和尊崇归于他及圣父及圣神，从今世直到永远，世世无尽。阿们。\par

\SourceRecord{Translated by John A. Broadus.；Nicene and Post-Nicene Fathers, First Series；Vol. 13.；Edited by Philip Schaff.；Buffalo, NY: Christian Literature Publishing Co.,；1889.；<http://www.newadvent.org/fathers/23055.htm>.}
